\documentclass[tikz,border=2mm]{standalone}
\usetikzlibrary{math}

\makeatletter

\newif\iffounddigit

% === Improved scanner: expand first, then look for digits ===
\def\scan@for@digits#1{%
  \ifx#1\@empty
  \else
    \ifnum`#1>47 \ifnum`#1<58 \global\founddigittrue\fi\fi
    \expandafter\scan@for@digits
  \fi}

\newcommand{\checkfornumbers}[1]{%
  \global\founddigitfalse
  % Fully expand once, then detokenize
  \edef\temp@str{\detokenize{\expanded{#1}}}%
  \expandafter\scan@for@digits\temp@str\@empty
}

\tikzset{
  define/@circle-point/.code args={#1,#2}{%
    \message{^^J[PASS] Point: #1, #2^^J}%
    % do whatever you need for a center+point circle
  },
  define/@circle-radius/.code args={#1,#2}{%
    \message{^^J[MATCH] Radius: #1, #2^^J}%
    % do whatever you need for a center+radius circle
  },

  define/circle/.code args={#1,#2}{%
    \checkfornumbers{#2}%
    \iffounddigit
      \pgfkeys{/tikz/define/@circle-radius={#1,#2}}%
    \else
      \pgfkeys{/tikz/define/@circle-point={#1,#2}}%
    \fi
  },
}

\makeatother

\begin{document}
\begin{tikzpicture}
  \tikzmath{
    \O{1} = 1; \O{2} = 1;
    \P{1} = 2; \P{2} = 3;
    \radius = 5;
  }

  \tikzset{define/circle={\O,\P}}           % → Point
  \tikzset{define/circle={\O,\P{1}}}        % → Radius (has digit)
  \tikzset{define/circle={\O, 5cm}}         % → Radius
  \tikzset{define/circle={\O, \radius}}     % → Radius (now works!)
  \tikzset{define/circle={\O, 3.14}}        % → Radius
  \tikzset{define/circle={\O, abc}}         % → Point

\end{tikzpicture}
\end{document}